\documentclass[10pt,journal,compsoc]{IEEEtran}
\usepackage[utf8]{inputenc}
\usepackage[T1]{fontenc}
\usepackage{url}
\usepackage[hidelinks]{hyperref}
\usepackage{cite}
\usepackage{amsmath,amssymb,amsfonts}
\usepackage{algorithmic}
\usepackage{graphicx}
\usepackage{textcomp}
\usepackage{xcolor}
\usepackage{tabularx}
\usepackage{listings}
\usepackage{array}

\lstset{
  basicstyle=\small\ttfamily,
  breaklines=true,
  frame=single,
  columns=fullflexible
}

\begin{document}

\title{Module 005 Holographic Indexing}
\author{Bryan Alexander Buchorn \\ Independent Researcher \\ bryan.alexander@buchorn.com}
\markboth{CEREBRUM v1.2.4 Hardened Edition · March 2026}{}

\maketitle

\begin{abstract}
This module forms a critical component of the CEREBRUM framework, a community-structured graph attention engine for Knowledge Graph reasoning. It focuses on Advanced Graph Attention Analysis.
\end{abstract}

\begin{IEEEkeywords}
Knowledge Graph, Graph Attention, DSCF, CSA, Hierarchical Reasoning.
\end{IEEEkeywords}

\section{Holographic Indexing: Privacy-Preserving Discovery in Federated Knowledge Networks}

\textbf{Authors}: Bryan Alexander Buchorn  
\textbf{Affili\-ations}: Independent Researcher, Las Vegas, NV, USA  
\textbf{Status}: v1.2.0 (Hardened Enterprise)  
\textbf{Date}: March 2026

---

\subsubsection{Abstract}
Federated Knowledge Graph reasoning requires nodes to identify relevant information across decen\-tralized peers without compr\-omising data privacy or bandwidth. We present \textbf{Holographic Indexing}, a two-tier discovery protocol designed for "Blind Semantic Search." Our method combines \textbf{Bloom Filter} \cite{bloom1970} membership probing for exact entity matching with \textbf{Community Centroid Signatures} for semantic neigh\-borhood appro\-ximation. This "Holographic" repre\-sentation allows nodes to identify potential reasoning paths in remote peers with minimal data leakage. We formalize the const\-ruction of these signatures and the \textbf{Wormhole Attention} mechanism that enables cross-graph traversals. We further describe security enhan\-cements in v1.2.0, including \textbf{HMAC-SHA256 Path Provenance} \cite{gentry2009} and \textbf{Federated Leases}, which ensure the integrity and reliability of federated reasoning in high-velocity, multi-tenant envir\-onments.

\subsubsection{1. Introd\-uction}
The expansion of decen\-tralized Knowledge Graphs neces\-sitates a protocol for inter-graph discovery that respects the data sovereignty of individual nodes. Traditional federated search methods often require a central index or the exchange of full node lists, both of which are unacc\-eptable in privacy-sensitive domains (e.g., healthcare or defense). Holographic Indexing addresses this by exchanging compressed, non-reversible topological signatures.

\subsubsection{2. Tier 1: The Entity Hologram}
Each node $\mathcal{G}_i$ generates a Bloom Filter $B_i$ of its entity set $\mathcal{E}_i$. This allows a peer $\mathcal{G}_j$ to verify the existence of a specific entity $e$ with a confi\-gurable false-positive rate $p$, while ensuring that the full set $\mathcal{E}_i$ cannot be enumerated via the filter.

\subsubsection{3. Tier 2: The Community Hologram}
To support fuzzy, semantic discovery, we introduce the Community Centroid Signature. For every community $C_k \in \mathcal{G}_i$, the node computes:
\begin{itemize}
\item \textbf{Centroid}: $\vec{c}_k = \text{mean}(\{\vec{e} : e \in C_k\})$
\item \textbf{Radius}: $\rho_k = \max \|\vec{e} - \vec{c}_k\|$

\end{itemize}
The set of all centroids forms the "Semantic Hologram." A remote reasoning beam looking for concept $\vec{x}$ can compute the "Wormhole Score":
\begin{equation}
\text{score}(C_k) = \cos(\vec{x}, \vec{c}_k)
\end{equation}
Peers exceeding a threshold $\sigma$ (default 0.75) are flagged as relevant reasoning desti\-nations.

\subsubsection{4. Enterprise Security (v1.2.0)}
To prevent adversarial path injection in federated networks, v1.2.0 implements \textbf{HMAC-SHA256 Path Provenance}. Every reasoning response from a remote adapter is crypt\-ographically signed using a shared secret $\mathcal{K}$. This ensures that "Wormhole" paths are both seman\-tically valid and struc\-turally authentic.

\subsubsection{5. Conclusion}
Holographic Indexing provides a mathe\-matically robust and privacy-preserving framework for federated graph reasoning. By decoupling exact membership from semantic relevance, it enables secure, decen\-tralized intel\-ligence at scale.

---
\textbf{References}
\begin{enumerate}
\item Bloom, B. H. (1970). Space/time trade-offs in hash coding with allowable errors. Commun\-ications of the ACM.
\item Gentry, C. (2009). Fully homomorphic encryption using ideal lattices. STOC.
\item Kairouz, P., et al. (2021). Advances and open problems in federated learning. Foundations and Trends in Machine Learning.
\item Broder, A., \& Mitzen\-macher, M. (2004). Network appli\-cations of Bloom filters: A survey. Internet Mathematics.
\item Tarkoma, S., Rothenberg, C. E., \& Lagerspetz, E. (2011). Theory and practice of Bloom filters for distributed systems. IEEE Commun\-ications Surveys \& Tutorials.
\item Buchorn, B. A., \& Sonnet, C. (2026). Federated HMAC Security in CEREBRUM. SPEC\_005.md.

\end{enumerate}

\bibliographystyle{IEEEtran}
\bibliography{../../docs/latex/references}

\end{document}
